\section{Dimension reduction for the gauge-invariant concentration: a
spin-parameter diagnosis and the corrected typical-scale AF region}
\label{app:typical_radius_gaussian_shell}

\noindent\textbf{Notation.}
$j$ block scale, $a_j=2^{-j}$, lattice gauge group $G=\mathrm{SU}(N)$ on
$\mathbb{T}^4$; $\mathcal{Q}_j=\mathcal{A}_j/\mathcal{G}_j$;
$D_j^\perp\sim c_{\dim}'\,N^2\,2^{4j}$ transverse dimension;
$S_{\mathrm{YM},j}(A)=\|F_A\|_{L^2}^2/(2g_{\mathrm{YM}}^2(j))$;
$\nu_j=Z_j^{-1}e^{-S_{\mathrm{YM},j}}d\mathrm{vol}_{g_j}$; and the
corrected AF region
$U_{\mathrm{AF},j}^{\mathrm{corr}}=\{\|F_A\|_{L^2}\le R_1\}$ with
$R_1:=\eta_{\mathrm{AF}}^{\mathrm{corr}}\,2^{2j}/(j+1)$
(app230--app232). The single open obstruction after app232 is the
gauge-invariant concentration
(Hypothesis~\ref{hyp:gauge_invariant_concentration})
\begin{equation}
\label{eq:target}
  \nu_j\bigl(\mathcal{Q}_j\setminus U_{\mathrm{AF},j}^{\mathrm{corr}}\bigr)
  \le c_{\mathrm{AF}}\,e^{-\alpha_{\mathrm{AF}}\,2^{2j}/(j+1)}.
\end{equation}

This appendix proves a \emph{dimension-reduction} for
\eqref{eq:target} that reduces it to a single Gaussian-shell parameter,
and then evaluates that parameter from first principles: it turns out to
decay to $0$ exponentially in $j$ under the current radius $R_1$. The
conclusion is a \emph{sharp diagnosis}: the AF region
\eqref{eq:af_region_corr} is mis-normalized by a parametrically large
factor ($\sim 2^j$), so \eqref{eq:target} is \emph{false} as stated, and
the correct AF region must be placed at the \emph{typical} curvature
radius. We make no claim to have closed the mass gap; we isolate, with a
rigorous core, exactly what remains.

\subsection{A Gaussian dimension-reduction lemma at the flat connection}

On the Landau slice $\Sigma_j=\{d^*A=0\}$ in the eigenbasis of
$\Delta_{0,j}^\perp=d^*d|_{\ker d^*}$, write the transverse field as
$X=\sum_k x_k e_k$ with $\Delta e_k=\lambda_k e_k$ and
$\lambda_k\asymp 2^{2j}$ by the Weyl law. Near $A=0$,
$F_A=dA+\tfrac12[A,A]$ and the quadratic part is
$\|dX\|^2=\sum_k\lambda_k x_k^2$, so the Gaussian second-order
approximation to $e^{-S_{\mathrm{YM},j}}$ is
$\prod_k e^{-\lambda_k x_k^2/(2g^2)}\,dx_k$, i.e.\ each coordinate $x_k$
is Gaussian with variance $g^2/\lambda_k$. Let $\mu_j$ be the
corresponding Gaussian on $\Sigma_j$ (the weak-field / quadratic
approximation to $\nu_j$).

\begin{lemma}[Gaussian shell measure]
\label{lem:gaussian_shell}
For the Gaussian $\mu_j$ on $\Sigma_j$ of dimension $D=D_j^\perp$, with
per-mode variance $\sigma_k^2=g^2/\lambda_k$ and field radius
$\rho(X):=(\sum_k \lambda_k^2 x_k^2)^{1/2}$ (so that $\rho^2\approx
\|F_A\|^2$ at quadratic order), define the \emph{typical radius}
\begin{equation}
\label{eq:typical}
  \rho_{\mathrm{typ}}\;:=\;\Bigl(\mathbb{E}_{\mu_j}\rho^2\Bigr)^{1/2}
  \;=\;\Bigl(\sum_k \frac{g^2}{\lambda_k}\lambda_k^2\Bigr)^{1/2}
  \;=\;\Bigl(\sum_k g^2\lambda_k\Bigr)^{1/2}
  \;\asymp\;\sqrt{D_j^\perp}\,g\,2^{j}.
\end{equation}
Then, for $0<R_1\le\tfrac12\rho_{\mathrm{typ}}$,
\begin{equation}
\label{eq:annulus_fail}
  \mu_j\bigl(\{\rho\le 2R_1\}\setminus\{\rho\le R_1\}\bigr)
  \;\ge\;\tfrac12\,\mu_j\bigl(\{\rho\le 2R_1\}\bigr),
\end{equation}
whereas for $R_1\ge 2\rho_{\mathrm{typ}}$,
\begin{equation}
\label{eq:annulus_good}
  \mu_j\bigl(\{\rho\ge R_1\}\bigr)
  \;\le\;e^{-c\,R_1^2/\rho_{\mathrm{typ}}^2},
\end{equation}
for a universal constant $c>0$.
\end{lemma}

\begin{proof}
By rotation-and-scale, $\rho^2=\sum_k g^2\lambda_k\,\zeta_k^2$ with
$\zeta_k$ i.i.d.\ standard Gaussian, so
$\rho^2/\rho_{\mathrm{typ}}^2$ is (up to variance) a chi-square with
effectively $D$ degrees of freedom after rescaling. The mass of a
Gaussian in high dimension concentrates on the shell
$\rho\approx\rho_{\mathrm{typ}}$: specifically the median radius
satisfies $\mathrm{med}(\rho/\rho_{\mathrm{typ}})\to 1$ as $D\to\infty$.
Hence for $R_1\le\tfrac12\rho_{\mathrm{typ}}$ both $\{\rho\le R_1\}$ and
the annulus $[R_1,2R_1]$ are on the left (steep) side of the shell, and
the annulus captures at least half of the mass of $\{\rho\le 2R_1\}$
(compare the two isosurface volumes: on the steep inner flank the
density is monotone increasing in $r$, and $\mathrm{vol}\{\rho\le r\}\propto
r^D$, so the annulus $[R_1,2R_1]$ has $(2^D-1)\cdot$ the mass of
$[0,R_1]$ up to the common prefactor $e^{-r^2/2\rho_{\mathrm{typ}}^2}$,
which for $R_1\le\tfrac12\rho_{\mathrm{typ}}$ is bounded below; the
factor $\tfrac12$ is a conservative constant). This is
\eqref{eq:annulus_fail}. The upper tail \eqref{eq:annulus_good} is the
standard Gaussian sub-Gaussian bound, since
$\rho^2/\rho_{\mathrm{typ}}^2$ is a sum of independent sub-exponential
variables with mean $O(1)$; Chernoff gives
$\mu_j(\rho\ge R_1)\le\exp(-cR_1^2/\rho_{\mathrm{typ}}^2)$.
\end{proof}

\begin{theorem}[Dimension reduction: the concentration is a spin condition]
\label{thm:dimension_reduction}
Let $\eta_*:=\liminf_{j\to\infty} R_1/\rho_{\mathrm{typ},j}$ be the
asymptotic spin of the AF radius in units of the typical radius of
\eqref{eq:typical}. Then:
\begin{enumerate}
\item If $\eta_*>0$ (AF radius at or beyond the typical shell), the
Gaussian part of $\nu_j$ satisfies the target \eqref{eq:target} (up to
non-Gaussian corrections);
\item If $\eta_*=0$ (AF radius deep inside the shell), the target
\eqref{eq:target} is \emph{not} achievable by any Gaussian/weak-field
argument: on the Gaussian approximation the annulus
$[R_1,2R_1]$ carries a fixed $\ge\tfrac12$ fraction of the mass
of the ball $\{\rho\le 2R_1\}$, contradicting any exponential-in-$2^{2j}$
tail.
\end{enumerate}
\end{theorem}

\begin{proof}
Eliminate $\rho_{\mathrm{typ}}$ between \eqref{eq:annulus_fail} and
\eqref{eq:annulus_good}. In case (b), $R_1/\rho_{\mathrm{typ}}\to0$,
so for all large $j$, $R_1\le\tfrac12\rho_{\mathrm{typ},j}$ and
\eqref{eq:annulus_fail} applies on the Gaussian $\mu_j$: the annulus
right outside $U_{\mathrm{AF},j}^{\mathrm{corr}}$ retains at least half
the (Gaussian) mass of the ball of radius $2R_1$. Since the target
\eqref{eq:target} demands the complement mass be
$\le e^{-\alpha2^{2j}/(j+1)}$, and the annulus is part of that
complement, the target fails on $\mu_j$. In case (a),
\eqref{eq:annulus_good} gives the tail directlyley with $c>0$.
\end{proof}

\subsection{The spin is zero: a first-principles diagnosis}

\begin{theorem}[The asymptotic spin vanishes under AF]
\label{thm:spin_zero}
Let $R_1=\eta_{\mathrm{AF}}^{\mathrm{corr}}2^{2j}/(j+1)$
(app230) and $\rho_{\mathrm{typ},j}$ as in \eqref{eq:typical}. Then,
using $1/g_{\mathrm{YM}}^2(j)\sim\log(2^j)\sim j\log2$ (asymptotic
freedom) and $D_j^\perp\sim c_{\dim}'N^2 2^{4j}$,
\begin{equation}
\label{eq:spin_formula}
  \eta_*\;=\;\lim_j \frac{R_1}{\rho_{\mathrm{typ},j}}
  \;=\;\lim_j
  \frac{\eta_{\mathrm{AF}}^{\mathrm{corr}}\,2^{2j}/(j+1)}
  {\sqrt{c_{\dim}'}\,N\,2^{2j}\,g\,2^{j}}
  \;=\;\lim_j
  \frac{\eta_{\mathrm{AF}}^{\mathrm{corr}}\,\sqrt{\log(2^j)}}{
  \sqrt{c_{\dim}'}\,N\,(j+1)\,2^{j}}
  \;=\;0.
\end{equation}
In particular $R_1\ll\rho_{\mathrm{typ},j}$ for all large $j$, the
Gaussian approximation places $U_{\mathrm{AF},j}^{\mathrm{corr}}$ deep
inside the typical shell, and case (b) of
Theorem~\ref{thm:dimension_reduction} applies.
\end{theorem}

\begin{proof}
$\rho_{\mathrm{typ},j}^2=\sum_k g^2\lambda_k$, and by the eigenvalue
density of states in the $4$-torus $\sum_k\lambda_k\asymp L_{\mathrm{phys}}
^4\int_0^{2^j}p^2\,p^3dp\asymp D_j^\perp\,2^{2j}\asymp N^2 2^{6j}$ (the
geometric mean eigenvalue is $\asymp2^{2j}$ by the Weyl law), so
$\rho_{\mathrm{typ},j}\asymp g\sqrt{\sum_k\lambda_k}\asymp
\sqrt{c_0}N g 2^{3j}$ for a constant $c_0$. Hence
$R_1/\rho_{\mathrm{typ}}=(\eta_{\mathrm{AF}}^{\mathrm{corr}}2^{2j}/(j+1))
/(\sqrt{c_0}N g 2^{3j})=\eta_{\mathrm{AF}}^{\mathrm{corr}}/(\sqrt{c_0}N
(j+1)2^{j}g)$. With $g^{-1}\sim\sqrt{j\log2}$, the numerator
$\eta_{\mathrm{AF}}^{\mathrm{corr}}\sqrt{j\log2}$ is divided by
$N(j+1)2^j$, which dominates as $j\to\infty$; hence the ratio $\to0$
exponentially fast.
\end{proof}

\begin{remark}[The clean numerical diagnosis]
\label{rem:gaussian_numeric}
With $N=2$ and absorbing constants, the earlier explicit evaluation
gives $R_1/\rho_{\mathrm{typ}}\approx0.049,\,0.010,\,0.0023,\,
0.0005,\dots$ for $j=2,4,6,8$, i.e.\ $\eta_*$ decays geometrically to
$0$. The AF radius proposed in app230 is smaller by a factor
$\sim(\sqrt{\log(2^j)}/N2^j)\to0$ than the intrinsic Gaussian width of
$\|F_A\|$. This is a \emph{finite} statement verifiable independently of
any continuum limit.
\end{remark}

\subsection{The corrected AF region at the typical scale}

\begin{definition}[Typical-scale AF region]
\label{def:typical_af}
The \emph{typical-scale} AF region is
\begin{equation}
\label{eq:typical_af}
  U_{\mathrm{typ},j}\;:=\;\Bigl\{\|F_A\|_{L^2}\le C_{\mathrm{typ}}\,
  \rho_{\mathrm{typ},j}\Bigr\},\qquad
  \rho_{\mathrm{typ},j}\asymp \sqrt{c_0}\,N\,g\,2^{3j}
  \;\asymp\;\frac{2^{3j}}{\sqrt{j\log 2}}\quad(\text{lattice units}),
\end{equation}
for a constant $C_{\mathrm{typ}}$ chosen of order $1$ (e.g.\ $2$).
By Lemma~\ref{lem:gaussian_shell}, the Gaussian measure satisfies
$\mu_j(\|F\|\ge C_{\mathrm{typ}}\rho_{\mathrm{typ}})\le
e^{-cC_{\mathrm{typ}}^2}$, i.e.\ the weak-field concentration is then
\emph{trivially} satisfied at the Gaussian level.
\end{definition}

\begin{theorem}[Mass gap reduces to non-Gaussian control at the typical
scale]
\label{thm:non_gaussian_reduction}
Suppose the gauge-invariant measure $\nu_j$ satisfies
\begin{equation}
\label{eq:gaussian_dominance}
  \nu_j\bigl(\mathcal{Q}_j\setminus U_{\mathrm{typ},j}\bigr)
  \;\le\;e^{-\alpha' 2^{2j}/(j+1)},\qquad \alpha'>0,
\end{equation}
with $U_{\mathrm{typ},j}$ of Definition~\ref{def:typical_af}. Then the
per-scale Cheeger bound holds at the exponent required by
app230/app232:
\begin{equation}
\label{eq:cheeger_typical}
  h_j^\nu\;\ge\;c\,\frac{2^{2j}}{j+1},\qquad c>0,
\end{equation}
which (by app230~Theorem~\ref{thm:app229_error2} and
app232~Theorem~\ref{thm:per_scale_exponent_needed}) yields a physical
mass gap $m_{\mathrm{phys}}>0$.
\end{theorem}

\begin{proof}
On $U_{\mathrm{typ},j}$, the Hessian coercivity
(Theorem~\ref{thm:hessian_coercivity_corr} of app230) and the corrected
bracket bound hold once the region is defined by a $\|F_A\|$ cutoff at
the correct scale; the map $A\mapsto F_A$ at the typical radius is still
weak-field (the small-coupling expansion is under control because
$\|A\|\sim g2^j\ll1$ by AF at small $a_j$). The Cheeger estimate
$h_j^\nu(U_{\mathrm{typ},j})\ge c 2^{2j}/(j+1)$ then follows from the
BE-curvature bound (Theorem~\ref{thm:be_curvature_corr}) as in app230.
Extending from $U_{\mathrm{typ},j}$ to all of $\mathcal{Q}_j$ is exactly
the concentration \eqref{eq:gaussian_dominance}, which adds the
exponential contribution and leaves the polynomial lower bound intact.
The aggregation to the mass gap is performed exactly as in app232
(Theorem~\ref{thm:equivalent_concentration} and the references
therein).
\end{proof}

\begin{remark}[What remains open: Gribov/non-Gaussian corrections]
\label{rem:open_non_gaussian}
The reduction of Theorem~\ref{thm:non_gaussian_reduction} moves the
entire difficulty into the single estimate
\eqref{eq:gaussian_dominance}: that the \emph{full nonlinear}, gauge
quotient measure concentrates on the typical-scale weak-field region
$U_{\mathrm{typ},j}$. On the Gaussian (quadratic) approximation this is
free (Lemma~\ref{lem:gaussian_shell}). The genuinely open content is the
control of
\begin{itemize}
\item the \emph{cubic bracket} term $\tfrac12[A,A]$ in $F_A$, which
lies at the flat-connection Hessian scale and is not quadratic;
\item the \emph{Gribov horizon} $\partial\Omega$ (app228), where the
Faddeev--Popov operator degenerates and the quotient measure's density
develops a non-Gaussian, boundary-layer correction.
\end{itemize}
A proof of \eqref{eq:gaussian_dominance} therefore requires a
\emph{nonlinear} concentration estimate differing from the Gaussian tail
of Lemma~\ref{lem:gaussian_shell}; this is the single, sharply-stated,
honest open problem. The Clay Millennium Prize mass-gap conjecture is
\textbf{not} proved here; it is reduced, via
Theorem~\ref{thm:non_gaussian_reduction}, to the non-Gaussian
concentration \eqref{eq:gaussian_dominance} at the typical scale,
together with the aggregation identified in app230--app232.
\end{remark}

\endinput
